\documentclass[10pt]{article}
\usepackage[utf8]{inputenc}
\usepackage[brazil]{babel}
\usepackage[paperwidth=11cm,paperheight=6.4cm,margin=.55cm]{geometry}
\usepackage[sfdefault]{roboto}
\usepackage{xcolor,booktabs,tabularx}
\definecolor{UFGBlue}{HTML}{0072B9}
\pagestyle{empty}
\setlength{\parindent}{0pt}
\setlength{\parskip}{3pt}
\begin{document}
{\small\color{UFGBlue} TRECHO DO ARQUIVO DO PROJETO}\par
{\large\bfseries Especificação de requisitos}\par
\textbf{RF21 Registrar o contexto do feedback — essencial, E01.} O registro de feedback deve conter o fato motivador, os impactos observados, a data, o emissor, o receptor, a classificação (RF19) e o conteúdo, podendo associar um questionário. O sistema não verifica a veracidade dos fatos (RN16). O formulário deve orientar o emissor a não incluir dados sensíveis (RF-DOC-14).
\vfill
{\scriptsize especificacao-requisitos.md | seção 3.2}
\newpage
{\small\color{UFGBlue} TRECHO DO ARQUIVO DO PROJETO}\par
\begingroup
\fontsize{9}{10.5}\selectfont
\setlength{\parskip}{1.5pt}
{\large\bfseries CU-02 | Responder e enviar avaliação}\par
\textbf{Autoria:} Luiz Augusto Vaz.\par
\textbf{Ator principal:} Participante no papel de avaliador.\par
\textbf{Fluxo principal de sucesso}\par
1. O participante solicita ver suas avaliações pendentes.\par
2. O sistema apresenta as avaliações em rascunho das coletas abertas, com receptor, coleta, prazo e progresso de preenchimento.\par
3. O participante escolhe uma avaliação.\par
\vfill{\scriptsize CU-02-responder-e-enviar-avaliacao.md | recorte do fluxo}
\endgroup
\newpage
{\small\color{UFGBlue} TRECHO DO ARQUIVO DO PROJETO}\par
{\large\bfseries Matriz de rastreabilidade}\par
{\small Colunas selecionadas da matriz vertical.}\par
\small\renewcommand{\arraystretch}{1.4}\begin{tabularx}{\linewidth}{@{}llXll@{}}\toprule ID & Fonte & Caso & Aceitação & Situação\\\midrule
RF11 & E01 & CU-02 & CA01, CA07 & Proposto\\
RF12 & E01 & CU-02 & CA08 & Proposto\\
\bottomrule\end{tabularx}\par\vfill{\scriptsize matriz-rastreabilidade.md | seção 3.1}
\end{document}
